\documentclass[11pt,a4paper]{report}

% ---------- Codificación e idioma ----------
\usepackage[utf8]{inputenc}
\usepackage[T1]{fontenc}
\usepackage[spanish,es-noindentfirst]{babel}

% ---------- Diseño de página ----------
\usepackage[a4paper,margin=2.5cm,headheight=15pt]{geometry}
\usepackage{xcolor}
\usepackage{graphicx}
\usepackage{fancyhdr}
\usepackage{titlesec}
\usepackage{enumitem}
\usepackage[most]{tcolorbox}
\usepackage{booktabs}
\usepackage{tabularx}
\usepackage{caption}
\usepackage{needspace}
\usepackage{etoolbox}
\usepackage{amssymb}
\usepackage{hyperref}

% ---------- Paleta institucional (la misma de la app: docs/estilo-visual.md) ----------
% Morados del degradado de la barra superior y dorado de la pestaña activa.
\definecolor{umagpurple}{HTML}{3B28A3}
\definecolor{umagpurpledark}{HTML}{2D1B69}
\definecolor{umaggold}{HTML}{B08D57}
\definecolor{umaggray}{HTML}{4B5563}
\definecolor{cajafondo}{HTML}{F8FAFC}

\hypersetup{
  colorlinks=true,
  linkcolor=umagpurple,
  urlcolor=umagpurple,
  citecolor=umagpurple,
  pdftitle={Manual de Usuario - Sistema de Biblioteca UMAG},
  pdfauthor={Biblioteca UMAG},
}

% ---------- Encabezados / pies de página ----------
\pagestyle{fancy}
\fancyhf{}
\renewcommand{\headrulewidth}{0.4pt}
\renewcommand{\footrulewidth}{0pt}
% Encabezado de capítulo sin el prefijo "Capítulo N." (ya se ve grande en la
% portadilla del capítulo) para que quepa junto al nombre del manual sin superponerse.
\renewcommand{\chaptermark}[1]{\markboth{#1}{}}
\fancyhead[L]{\small\textcolor{umaggray}{Biblioteca UMAG}}
\fancyhead[R]{\small\textcolor{umaggray}{\leftmark}}
\fancyfoot[C]{\small\textcolor{umaggray}{\thepage}}

\fancypagestyle{plain}{%
  \fancyhf{}
  \fancyfoot[C]{\small\textcolor{umaggray}{\thepage}}
  \renewcommand{\headrulewidth}{0pt}
}

% ---------- Estilo de títulos ----------
\titleformat{\chapter}[display]
  {\normalfont\bfseries\color{umagpurple}}
  {\Large\filright\textcolor{umaggold}{CAPÍTULO \thechapter}}
  {6pt}
  {\Huge\filright}
\titlespacing*{\chapter}{0pt}{0pt}{28pt}

\titleformat{\section}
  {\normalfont\Large\bfseries\color{umagpurpledark}}{\thesection}{0.6em}{}
\titleformat{\subsection}
  {\normalfont\large\bfseries\color{umagpurple}}{\thesubsection}{0.6em}{}
% Un título de sección no queda solo al pie de la página: si no caben él,
% su etiqueta de permiso y unas líneas de texto, pasa a la página siguiente.
\pretocmd{\section}{\needspace{7\baselineskip}}{}{}

% ---------- Cajas reutilizables ----------
\newtcolorbox{nota}{
  colback=umaggold!8, colframe=umaggold, boxrule=0.6pt, arc=2pt,
  left=8pt, right=8pt, top=6pt, bottom=6pt,
  fonttitle=\bfseries, title={Tip},
}
\newtcolorbox{importante}{
  colback=umagpurple!6, colframe=umagpurple, boxrule=0.6pt, arc=2pt,
  left=8pt, right=8pt, top=6pt, bottom=6pt,
  fonttitle=\bfseries, title={Importante},
}
% Etiqueta compacta (no una caja que envuelve contenido) para marcar una
% sección que solo ven los roles con cierto permiso, sin arriesgar que una
% tcolorbox con una captura de pantalla adentro desborde el alto de la página.
% \requierepermiso{nombre del permiso, tal como aparece en Personal y Roles}
\newtcbox{\etiquetapermiso}{on line, colback=umagpurple, colframe=umagpurple,
  coltext=white, boxrule=0pt, arc=0pt, left=4pt, right=4pt, top=2pt, bottom=2pt,
  fontupper=\bfseries\small}
\newcommand{\requierepermiso}[1]{%
  \needspace{3\baselineskip}%
  \par\noindent\etiquetapermiso{REQUIERE PERMISO: #1}\par\vspace{6pt}%
}

% ---------- Comando para insertar capturas de pantalla ----------
% \captura{archivo}{texto de pie de figura}
\newcommand{\captura}[2]{%
  \begin{center}
  \fbox{\includegraphics[width=0.92\textwidth,height=0.78\textheight,keepaspectratio]{imagenes/#1}}\\[5pt]
  {\small\itshape\color{umaggray} #2}
  \end{center}
  \vspace{6pt}
}

\setlist[itemize]{leftmargin=1.4em}
\setlist[enumerate]{leftmargin=1.6em}

\begin{document}

% ======================================================================
% PORTADA
% ======================================================================
\begin{titlepage}
\centering
\vspace*{1.5cm}

{\includegraphics[width=3cm]{imagenes/logo-umag.png}}

\vspace{0.8cm}

{\color{umaggray}\large UNIVERSIDAD DE MAGALLANES}

\vspace{2.2cm}

{\color{umagpurple}\Huge\bfseries Sistema de Biblioteca}\\[6pt]
{\color{umagpurpledark}\Huge\bfseries UMAG}

\vspace{1cm}
{\color{umaggold}\rule{8cm}{1.2pt}}
\vspace{1cm}

{\color{umaggray}\LARGE Manual de Usuario}

\vspace{0.4cm}
{\color{umaggray}\large Guía completa de funcionalidades para personal de biblioteca\\y usuarios del portal virtual}

\vfill

\begin{tcolorbox}[colback=cajafondo,colframe=umagpurple!40,width=0.8\textwidth,arc=3pt,boxrule=0.6pt]
\centering
\textbf{Panel de personal:} ruta \texttt{/login}\\[2pt]
\textbf{Portal de usuarios:} ruta \texttt{/portal/login}\\[6pt]
{\small (accede desde la dirección web que te indique el equipo de sistemas de tu institución)}
\end{tcolorbox}

\vspace{1cm}
{\color{umaggray} Versión del documento: septiembre de 2026}

\vspace{1cm}
\end{titlepage}

\tableofcontents
\clearpage

% ======================================================================
\chapter{Introducción}
% ======================================================================

\section{Sobre este manual}

Este documento explica, paso a paso, cómo usar el Sistema de Biblioteca UMAG.
Cubre las dos formas de acceder al sistema:

\begin{itemize}
  \item El \textbf{panel de personal}, usado por quienes trabajan en la
  biblioteca (asistentes de mesón, bibliotecarios y administradores) para
  atender préstamos, registrar entradas, gestionar el catálogo, generar
  reportes, etc.
  \item El \textbf{portal del usuario}, un sitio de autoservicio donde
  estudiantes, docentes y funcionarios pueden marcar su propia entrada,
  consultar el catálogo, reservar salas de estudio y descargar su
  constancia de no multa, sin necesidad de pasar por mesón.
\end{itemize}

Cada capítulo corresponde a una pantalla o funcionalidad del sistema,
acompañada de capturas de pantalla reales y una explicación de qué hace
cada botón, campo o filtro.

\begin{nota}
Las secciones marcadas con el recuadro \textbf{``Requiere permiso''} solo
aparecen para las cuentas cuyo rol incluye ese permiso. Qué incluye cada
rol se puede ajustar (sección~\ref{cap:personal-roles}); los roles que
trae el sistema de fábrica se resumen en la
sección~\ref{sec:roles}.
\end{nota}

\section{Requisitos}

\begin{itemize}
  \item Un navegador web moderno (Chrome, Edge, Firefox).
  \item Conexión a la red donde esté publicado el sistema.
  \item Credenciales de acceso: un correo y contraseña de personal, o un
  RUT y contraseña de usuario (estudiante/docente/funcionario), según
  corresponda.
\end{itemize}

\section{Dos sistemas de acceso independientes}

Es importante entender que el sistema tiene \textbf{dos puertas de entrada
completamente separadas}, cada una con su propio inicio de sesión:

\begin{enumerate}
  \item El \textbf{panel de personal} (\texttt{/login}), pensado para
  quienes trabajan en la biblioteca.
  \item El \textbf{portal del usuario} (\texttt{/portal/login}), pensado
  para estudiantes, docentes y funcionarios.
\end{enumerate}

Una cuenta de personal no sirve para entrar al portal de usuarios, y
viceversa. Si alguien te dice que ``no puede entrar'', lo primero que hay
que verificar es que esté usando la puerta correcta.

% ======================================================================
\chapter{Acceso para el Personal de Biblioteca}
% ======================================================================

\section{Inicio de sesión}

Para ingresar al panel de personal, abre la dirección del sistema y
completa el formulario con tu correo institucional y contraseña.

\captura{01-login-staff.png}{Pantalla de inicio de sesión del panel de personal.}

Si las credenciales son incorrectas, o si tu cuenta fue desactivada por un
administrador, el sistema mostrará un mensaje de error y no dejará
continuar.

\begin{nota}
Existe una segunda variante de esta pantalla (``Versión 2''), enlazada
desde la misma página de login, con un diseño alternativo. Cumple
exactamente la misma función; es solo una opción visual distinta.
\end{nota}

\section{Roles del personal}
\label{sec:roles}

Cada persona que atiende tiene \textbf{su propia cuenta}, y cada cuenta
tiene un \textbf{rol}. El rol define qué módulos aparecen en el menú y qué
botones se pueden usar. El sistema trae tres roles de fábrica:

\begin{itemize}
  \item \textbf{Atención de público:} el personal de mesón. Registra
  entradas, presta y recibe libros y equipos, y reserva salas. Puede
  consultar usuarios (y habilitar a quien esté inactivo), multas y
  listados, pero no cataloga, no cambia estados de libros ni la
  configuración.
  \item \textbf{Bibliotecario/a:} lo anterior, más los procesos técnicos y
  la supervisión: catalogación, estados de libros (incluida la baja),
  cambio masivo, equipos, cobro y anulación de multas, historiales y
  reportes.
  \item \textbf{Administrador/a:} acceso total, incluidos el Registro
  Absoluto, la configuración institucional, los catálogos del sistema,
  desactivar usuarios y las cuentas y roles del personal.
\end{itemize}

\begin{center}
\small
\begin{tabularx}{\textwidth}{@{}>{\raggedright\arraybackslash}Xccc@{}}
\toprule
\textbf{Qué se puede hacer} & \textbf{\shortstack{Atención\\de público}} & \textbf{Bibliotecario/a} & \textbf{Administrador/a} \\
\midrule
Registrar entradas & \checkmark & \checkmark & \checkmark \\
Préstamos, devoluciones y reservas de libros & \checkmark & \checkmark & \checkmark \\
Salas de estudio & \checkmark & \checkmark & \checkmark \\
Consultar usuarios y Constancia de No Multa & \checkmark & \checkmark & \checkmark \\
Habilitar usuarios inactivos & \checkmark & \checkmark & \checkmark \\
Ver multas & \checkmark & \checkmark & \checkmark \\
Listados de préstamos y de libros & \checkmark & \checkmark & \checkmark \\
Mostrar el código QR de acceso & \checkmark & \checkmark & \checkmark \\
Reservar salas en horas ya pasadas & --- & \checkmark & \checkmark \\
Registrar pagos, cursar y anular multas & --- & \checkmark & \checkmark \\
Catalogar libros e imprimir etiquetas de código & --- & \checkmark & \checkmark \\
Cambiar el estado de un libro (incluida la baja) & --- & \checkmark & \checkmark \\
Cambio masivo de estado & --- & \checkmark & \checkmark \\
Administrar equipos & --- & \checkmark & \checkmark \\
Historiales y reportes & --- & \checkmark & \checkmark \\
Regenerar el código QR & --- & \checkmark & \checkmark \\
Registro Absoluto & --- & --- & \checkmark \\
Configuración institucional y catálogos & --- & --- & \checkmark \\
Desactivar usuarios & --- & --- & \checkmark \\
Cuentas y roles del personal & --- & --- & \checkmark \\
\bottomrule
\end{tabularx}
\end{center}

Un administrador puede cambiar qué incluye cada rol o crear roles nuevos
(por ejemplo, ``Ayudante de sala'') desde \textbf{Personal y Roles}
(sección~\ref{cap:personal-roles}). Tu rol aparece en el Panel Principal
y bajo tu nombre, en la barra superior.

Si intentas abrir un módulo que tu rol no incluye (por ejemplo, escribiendo
su dirección a mano), el sistema te devuelve al Panel Principal con un
aviso.

\section{Cambiar mi contraseña}

En la barra superior, pulsa tu nombre (la burbuja con tus iniciales,
arriba a la derecha) y elige \textbf{Cambiar mi contraseña}: se pide la
actual y dos veces la nueva (mínimo 8 caracteres). Al pulsar \textbf{Guardar}, el
sistema pide confirmar el cambio; al confirmarlo se cierran tus otras
sesiones abiertas, por ejemplo en otro computador.

\captura{31-cambiar-contrasena.png}{Formulario para cambiar la propia contraseña, desde el menú de tu cuenta.}

Si olvidaste tu contraseña, pide a un administrador que la restablezca
desde Personal y Roles.

\begin{nota}
Las acciones que cambian datos importantes ---cerrar sesión, cambiar
contraseñas, roles o permisos, cursar multas, liberar o cancelar
reservas, dar de baja equipos o exportar a Excel--- piden confirmación
antes de aplicarse, con un resumen de lo que va a pasar. \textbf{Cancelar}
no cambia nada.
\end{nota}

\begin{nota}
Las búsquedas por título, autor o nombre no distinguen tildes, ñ ni
mayúsculas: ``calculo'' encuentra ``Cálculo I'' y ``munoz'', ``Muñoz''.
\end{nota}

\section{Guía ``¿Cómo se usa esta pantalla?''}

Las pantallas de uso diario tienen un botón con los pasos básicos de esa
pantalla, pensados para personal nuevo: \textbf{¿Cómo se usa?}, arriba a
la derecha en Entradas y Préstamos, o \textbf{¿Cómo se usa esta
pantalla?}, bajo el título, en las demás. Al pulsarlo se despliegan los
pasos; con \textbf{Ocultar} se vuelven a esconder. Cada vez que se entra a
la pantalla, la guía aparece cerrada.

\captura{32-guia-uso.png}{Guía de uso desplegada en la pantalla de Préstamos.}

% ======================================================================
\chapter{Panel Principal (Inicio)}
% ======================================================================

Al iniciar sesión, llegas al \textbf{Panel Principal}. Desde aquí puedes
ver de un vistazo cuántas personas hay en la biblioteca en este momento y
cuántas entradas se registraron hoy, además de accesos rápidos a los
módulos más usados.

\captura{02-dashboard.png}{Panel Principal, con indicadores del día, tu rol y accesos rápidos.}

\begin{itemize}
  \item \textbf{Personas en biblioteca:} cantidad de entradas registradas
  hoy (el sistema no distingue si la persona ya se retiró; es un conteo
  de ingresos del día, no de ocupación en tiempo real).
  \item \textbf{Entradas hoy:} total de registros de entrada del día.
  \item \textbf{Tu rol:} el rol de tu cuenta (sección~\ref{sec:roles}).
  \item \textbf{Accesos rápidos:} atajos a Registrar Entrada, Préstamos,
  Reservar Logia y Reportes, solo los que tu rol puede abrir. Cada
  tarjeta indica su atajo de teclado: \texttt{F1} Entradas, \texttt{F2}
  Préstamos, \texttt{F3} Salas y \texttt{F4} Reportes. Los atajos
  funcionan desde cualquier pantalla del panel.
\end{itemize}

\section{La barra superior}

La barra morada de arriba aparece en todas las pantallas del panel. De
izquierda a derecha tiene: el logo con la etiqueta \textbf{ADMIN} (vuelve
al Panel Principal); al centro, \textbf{Inicio} y los módulos principales
con su atajo de teclado; el botón \textbf{Gestión}
(capítulo~\ref{cap:gestion}), y tu burbuja, con tus iniciales, tu nombre
y tu rol. Al pulsarla se abre el menú de tu cuenta, con \textbf{Cambiar mi
contraseña} y \textbf{Cerrar sesión}, que pide confirmación antes de
salir. El módulo en el que estás aparece resaltado en dorado.

En pantallas angostas (celular o tablet vertical) los módulos se reúnen en
un solo botón \textbf{Menú}, que muestra el nombre del módulo actual y
abre el mismo menú de gestión con los módulos principales arriba.

% ======================================================================
\chapter{Registro de Entradas}
% ======================================================================

Este módulo permite marcar el ingreso de una persona a la biblioteca ---
ya sea un usuario institucional (estudiante, docente o funcionario), un
visitante externo o una visita puntual.

\requierepermiso{Registrar entradas}

\captura{03-entradas.png}{Registro de entrada: el RUT de quien entra, el resumen de hoy y las entradas del día.}

El botón \textbf{¿Cómo se usa?}, arriba a la derecha, despliega los pasos
para el personal nuevo.

\section{Cómo registrar una entrada, paso a paso}

Se escribe el RUT en la casilla de \textbf{Registro de entrada}. Los puntos y el guion se agregan solos mientras se
escribe (por ejemplo \texttt{15.003.151-6}).

\begin{nota}
Junto a la casilla hay un selector \textbf{RUT | Pasaporte / DNI}. Con
\textbf{RUT} (lo normal), el sistema revisa el dígito verificador antes de
buscar. Para un estudiante extranjero se marca \textbf{Pasaporte / DNI} y
se escribe el número tal como aparece en su documento, sin formato de RUT.
Si con RUT marcado se escribe una letra (que no sea la K final), el
selector pasa solo a Pasaporte / DNI. El mismo selector está en Préstamos,
Multas, las salas (uno por integrante) y en el portal.
\end{nota}

\captura{03-entradas-rut-lleno.png}{RUT escrito en el recuadro, listo para registrar.}

Al presionar \texttt{Enter} (o el botón \textbf{Registrar}), el
sistema valida que el RUT corresponda a un usuario activo y, si todo está
bien, confirma con un aviso verde y agrega la fila a las entradas de hoy
de inmediato --- sin recargar la página. La entrada queda a nombre de
quien la atendió, y el recuadro queda vacío y listo para la siguiente
persona.

\captura{03b-entradas-registrada.png}{Confirmación tras registrar: aviso verde con el nombre de la persona y la fila nueva en las entradas de hoy.}

\begin{nota}
Si el RUT no corresponde a ningún usuario, o el usuario está inactivo, el
sistema rechaza el registro y explica el motivo bajo el recuadro --- no hay
forma de ``forzar'' una entrada con un RUT inválido. Si la persona no es
de la universidad, se registra con el botón \textbf{Registrar visitante
externo}.
\end{nota}

\section{Tipos de entrada}

\begin{itemize}
  \item \textbf{Entrada normal (RUT):} para estudiantes, docentes y
  funcionarios ya registrados en el sistema. Se ingresa el RUT y el
  sistema valida que la persona exista y esté activa.
  \item \textbf{Visitante externo:} para quienes no pertenecen a la
  universidad. El botón \textbf{Registrar visitante externo} abre una
  ventana donde se elige el tipo de ingreso:
  \begin{itemize}
    \item \textbf{Externo:} una persona que no es de la universidad. Se
    registra su RUT y, si se quiere, su nombre, tal como los declara.
    \item \textbf{Visita:} una delegación o un colegio (por ejemplo, una
    comitiva o un tour). Queda marcada como ``Visita'' y se cuenta aparte
    en los reportes.
  \end{itemize}
\end{itemize}

\captura{03c-entradas-externo.png}{Ventana de Registrar visitante externo, con la opción Visita.}

\section{Hoy en la biblioteca}

El recuadro de la derecha resume el día de hoy, aunque abajo se esté
mirando otro:

\begin{itemize}
  \item Cuántas \textbf{personas distintas} marcaron su entrada hoy, frente
  al aforo de la biblioteca (220), con una barra de ocupación. Como el
  sistema no registra salidas, no es la cantidad de personas que están
  adentro en este momento.
  \item Las \textbf{entradas de hoy} ---una persona que sale y vuelve a
  entrar cuenta dos veces---, y cuántas fueron \textbf{por QR}, por el
  \textbf{portal digital} y \textbf{manuales}.
  \item Si la biblioteca está abierta («Abierta · cierra a las 20:30») o
  cuándo vuelve a abrir.
\end{itemize}

\section{Entradas del día}

Abajo están las entradas del día, de la más reciente a la más antigua, con
la hora, la persona, su RUT, su tipo (estudiante, docente, funcionario,
externo o visita) y la vía: \textbf{Manual} (registrada en mesón),
\textbf{Portal digital} (marcada por la persona con su RUT desde el
portal) o \textbf{QR} (escaneada por la persona desde su celular). Las flechas junto a
la fecha cambian de día ---al pulsar la fecha se elige otra en el
calendario--- y \textbf{Hoy} vuelve al día actual. Se muestran de a diez,
con páginas al pie.

\begin{itemize}
  \item El buscador filtra por nombre o RUT (con o sin puntos), sin
  importar las tildes.
  \item \textbf{Todas}, \textbf{Manual}, \textbf{QR} y \textbf{Externo}
  filtran según cómo entró la persona.
  \item La columna \textbf{Nota} dice la cuántas es esa entrada de la
  persona en el día. Si volvió a entrar a menos de diez minutos de la
  anterior, la nota sale en ámbar: puede ser un escaneo repetido.
  \item El botón de tres puntos de cada fila lleva a \textbf{Ver todas sus
  entradas} y, si es un usuario registrado, a \textbf{Ver sus préstamos}.
\end{itemize}

Para saber cuándo vino alguien en otros días, se escribe su nombre o RUT en
el buscador y se pulsa \textbf{Buscar en todas las fechas}: el listado pasa
a mostrar todo su historial, con la fecha de cada entrada, y se puede
acotar con \textbf{Desde} y \textbf{Hasta}. \textbf{Volver a hoy} regresa
al día actual.

\begin{nota}
Para ver todo lo que hizo una persona (entradas, préstamos y reservas
juntos) existe además el \textbf{Registro Absoluto}
(sección~\ref{cap:registro-absoluto}).
\end{nota}

% ======================================================================
\chapter{Préstamos, Devoluciones y Reservas}
% ======================================================================

Este es uno de los módulos más usados del sistema: permite prestar
libros y equipos (audífonos, notebooks, cargadores), registrar
devoluciones y reservar libros para retiro.

\requierepermiso{Préstamos y devoluciones}

\section{Buscar al usuario}

Todo comienza en el recuadro \textbf{RUT del usuario}. El RUT se puede
escribir con o sin puntos y guion (se ordena solo) o escanear desde el
carnet; luego se presiona \texttt{Enter} o \textbf{Buscar usuario}. Para un
estudiante extranjero se marca \textbf{Pasaporte / DNI} en el selector de
al lado. Antes
de buscar, la pantalla ya muestra lo que vendrá ---la ficha vacía, la tabla
de préstamos y los recuadros de equipos con cuántos hay disponibles---, todo
deshabilitado hasta elegir a alguien. El botón \textbf{¿Cómo se usa?} del
encabezado despliega los pasos para el personal nuevo.

\captura{04a-prestamos-buscar.png}{Pantalla de Préstamos antes de buscar: todo listo, esperando el RUT del usuario.}

Al encontrar al usuario, el sistema muestra:

\begin{itemize}
  \item Su ficha: nombre, RUT, tipo de usuario y correo; si su cuenta está
  \textbf{Habilitada} o \textbf{Bloqueada por multa}, y un resumen con sus
  préstamos activos, sus reservas y sus multas. Al lado, los botones
  \textbf{Nuevo préstamo}, \textbf{Reservar libro} y \textbf{Buscar otro
  usuario}.
  \item Un aviso rojo si tiene \textbf{multas pendientes}, con la cantidad
  y el monto total.
  \item Las pestañas \textbf{Préstamos} y \textbf{Reservas}, con lo vigente
  de cada una. En Préstamos están los libros y los equipos que tiene, con
  la fecha de devolución (y cuánto falta, o hace cuánto venció), el estado,
  la multa si la tiene y el botón \textbf{Devolver}. Lo ya devuelto queda
  guardado bajo \textbf{Ver devueltos}, salvo que tenga una multa por
  pagar. Al pie se indica cuántos libros tiene (el máximo es uno a la vez)
  y, si tiene reservas listas para entregar, un acceso directo a ellas.
\end{itemize}

\begin{importante}
Con \textbf{multas pendientes}, la persona \textbf{no puede llevarse ni
reservar material} (libros ni equipos): su ficha dice \textbf{Bloqueada por
multa}, y \textbf{Nuevo préstamo}, \textbf{Reservar libro} y los préstamos
de equipos quedan deshabilitados hasta que pague. Quien tenga el permiso
\textit{Registrar pagos} puede registrar el pago ahí mismo; al pagar o
anular la multa, se desbloquea al instante.
\end{importante}

\captura{04-prestamos.png}{Ficha de usuario con sus préstamos, tras buscar por RUT.}

\begin{nota}
Desde \textbf{Usuarios} (botón \textbf{Ver préstamos} de cada fila) y desde
\textbf{Entradas} (\textbf{Ver sus préstamos}) se llega a esta misma
pantalla con el RUT ya buscado.
\end{nota}

\section{Prestar un libro}

Con el botón \textbf{Nuevo préstamo} se abre el formulario para escanear
o escribir el código de barras del libro y elegir la fecha de devolución.
El sistema busca el ejemplar correspondiente y valida que esté disponible
y en estante antes de prestarlo.

\begin{nota}
La fecha de devolución tiene que caer en un día en que la biblioteca
atiende: si se elige un domingo, el formulario lo avisa en rojo y no deja
confirmar. Lo mismo vale para la fecha límite de retiro de una reserva.
\end{nota}

\begin{importante}
Cada usuario puede tener \textbf{un solo libro prestado a la vez}: si ya
tiene uno sin devolver, el sistema rechaza el segundo préstamo. Los
equipos no tienen ese límite.
\end{importante}

\begin{importante}
Si el libro que se quiere prestar ya está prestado o reservado por otra
persona, el sistema lo rechaza. En ese caso aparece la opción de
\textbf{unirse a la lista de espera} en su lugar --- cuando la copia se
devuelva, o cuando otra copia de ese título vuelva al estante (una nueva que
sale de inventario, o una que regresa de procesos técnicos), se aparta
automáticamente para la primera persona en la fila.
\end{importante}

\section{Reservar un libro}

Con el botón \textbf{Reservar libro} se aparta una copia para que el
usuario la retire más adelante, hasta la fecha límite que se elige. Si no
hay copias disponibles en ese momento, la persona entra a la cola de
espera igual que en un préstamo. Ambas cosas piden confirmación, igual que
\textbf{Cancelar} o \textbf{Salir de la fila} en la pestaña
\textbf{Reservas}. Cuando el plazo lo pone el sistema (al reservar desde
el portal o cuando le toca el turno a alguien de la lista de espera), es
de tres días; si cae en domingo, pasa al lunes. Si la persona no retira el
libro hasta la fecha límite, la reserva se cancela sola al comenzar el día
siguiente: la copia pasa a quien sigue en la lista de espera (que recibe el
aviso) o vuelve a quedar disponible, y queda anotado en la bitácora. Si
alguien con una reserva se lleva otra copia del mismo título, su reserva
queda como retirada con esa copia, y la que tenía apartada pasa a quien
sigue en la fila.

Cuando la persona viene a retirar su libro, se busca su RUT y, en la
pestaña \textbf{Reservas} (o con el acceso \textbf{reservas listas para
entregar} al pie de Préstamos), se pulsa \textbf{Entregar}: se abre
\textbf{Nuevo préstamo} con esa copia ya escaneada. Solo falta elegir la
fecha de devolución y confirmar; la reserva queda como retirada. Las
reservas ya retiradas o canceladas quedan guardadas bajo \textbf{Ver
retiradas y canceladas}.

\captura{04b-prestamos-entregar.png}{Tras pulsar Entregar: la copia reservada para esta persona, lista para prestar.}

\section{Préstamo de equipos}

Más abajo en la misma pantalla hay tres recuadros separados para prestar
\textbf{audífonos}, \textbf{notebooks} y \textbf{cargadores}; cada uno dice
cuántos quedan libres y, si la persona ya tiene uno, cuál. A diferencia de
los libros, estos préstamos no tienen fecha de devolución fija: se
devuelven al finalizar la estadía en la biblioteca. Al pulsar
\textbf{Pedir} se confirma el equipo (con su nombre de inventario) antes
de registrar el préstamo.

El equipo prestado aparece en la pestaña \textbf{Préstamos} con su
devolución \textbf{Pendiente} y se recibe con su botón \textbf{Devolver},
igual que un libro; desde ahí queda como \textbf{Devuelto}. Si no se
devolvió el mismo día, la fila lo marca \textbf{No se devolvió el mismo
día}.

\begin{nota}
El préstamo de un equipo se hace \textbf{escaneando su código de
barras físico}, igual que un libro --- no se tipea un nombre o código de
inventario a mano. El campo visible en pantalla (por ejemplo,
``Notebook 01'') es solo el nombre legible del equipo.
\end{nota}

\section{Devoluciones y multas por atraso}

Para recibir un libro o un equipo se pulsa \textbf{Devolver} en su fila y
se confirma.
Si el libro llega atrasado, el sistema \textbf{genera la multa solo}: una
tarifa fija por cada día de atraso (hoy \$15 por día, con un tope de
\$15.000 por préstamo), que queda pendiente a nombre del usuario y
registrada a nombre de quien recibió el libro. El aviso de confirmación
indica el monto.

Las cuentas con el permiso \textbf{Registrar pagos} ven además el botón
\textbf{Marcar pagada} junto a cada multa pendiente. El resto de la
gestión de multas (anular, cursar multas por daño o pérdida) está en el
módulo \textbf{Multas} (sección~\ref{cap:multas}).

\begin{nota}
Los equipos (audífonos, notebooks, cargadores) nunca generan multa, ya
que no tienen fecha de devolución fija.
\end{nota}

% ======================================================================
\chapter{Salas de Estudio (Logias)}
% ======================================================================

Este módulo administra la reserva de las 15 logias de estudio más las
salas con nombre propio (Sala de Seminarios, Sala de Postgrado y Sala
AGACI). Funciona en \textbf{horario continuo}, dentro del horario de
atención de la biblioteca: de lunes a viernes de 08:45 a 20:30 (la primera
reserva empieza a las 09:00, porque se agenda a la hora o a la media hora)
y el sábado de 09:00 a 14:00. Cada reserva dura entre 30 minutos y 2
horas. Lo que queda fuera de ese horario aparece rayado en la línea de
tiempo y no se reserva: el sábado, desde las 14:00, con un aviso arriba.
Los \textbf{domingos} la biblioteca no atiende: al elegir un domingo, la
línea de tiempo aparece rayada entera, con un aviso.

\requierepermiso{Salas de estudio}

\captura{05-salas.png}{Vista general de Salas de Estudio, con la línea de tiempo de cada sala.}

Cada sala tiene una línea de tiempo del día, con una línea guía por
cada hora. Los colores indican el estado de cada tramo: disponible
(verde, se puede hacer clic para reservar), por confirmar, plazo vencido,
confirmada o en uso, y llave devuelta. La línea azul marca la hora
actual ---arriba lleva la etiqueta \textbf{Ahora} con la hora--- y lo que
ya pasó del día aparece rayado.

La pantalla se actualiza sola cada 3 segundos mientras está a la vista,
y de inmediato al volver a ella: una reserva hecha desde el portal o
desde otro computador aparece sin recargar la página. Sobre la línea de
tiempo se indica a qué hora fue la última actualización, junto a un botón
para actualizar en el momento.

\section{Reservar una sala}

Se hace clic en el espacio verde de la sala, a la hora que se quiere, o
se usa \textbf{Reservar ahora} para una reserva que empieza de inmediato
(este botón aparece solo mientras la biblioteca atiende).
En el formulario se elige la hora de inicio, la duración, la cantidad de
personas (entre 2 y 5) y el RUT de cada una.

\captura{06-salas-modal-reserva.png}{Formulario de reserva, con los RUT del grupo ya completados y listo para confirmar.}

Si alguno de los RUT ingresados no corresponde a un usuario registrado,
o ya tiene otra reserva que choca en el horario, el sistema marca en
rojo exactamente el campo con el problema al intentar confirmar --- no
hace falta adivinar cuál de todos los RUT falló.

\begin{importante}
Todos los RUT ingresados deben corresponder a usuarios ya registrados en
el sistema --- no se pueden reservar logias para visitantes externos.
Además, el sistema no permite que la misma persona quede en dos salas
distintas al mismo tiempo, ni que sume más de \textbf{4 horas de reserva
en el mismo día}, contando todas las salas.
\end{importante}

\begin{nota}
Para el día de hoy solo se puede reservar desde la hora actual en
adelante. Registrar después una reserva que se atendió sin pasar por el
sistema (en una hora que ya pasó) requiere el permiso
\textbf{Reservas en horas ya pasadas}.
\end{nota}

\section{Confirmación de asistencia}

Cada grupo tiene \textbf{15 minutos desde la hora de inicio} de su reserva
(o desde el momento de reservar, si reservó con \textbf{Reservar ahora})
para presentarse y confirmar su llegada. Si el plazo vence sin
confirmación, la sala se libera automáticamente y queda disponible para
otra persona.

El panel \textbf{Menú de confirmación de asistencia}, en la parte
superior de esta pantalla, lista las reservas del día que todavía no
confirmaron llegada, con una cuenta regresiva en vivo. Desde ahí se puede:

\begin{itemize}
  \item \textbf{Confirmar llegada:} abre una ventana con los integrantes
  de la reserva, con su nombre y RUT, para verificar que estén antes de
  entregar la llave. La llegada se registra al pulsar \textbf{Confirmar}.
  \item \textbf{Liberar:} libera manualmente una reserva vencida, sin
  esperar a que expire sola. Pide confirmar, porque queda registrado que
  el grupo no se presentó.
\end{itemize}

\captura{06b-salas-confirmar-llegada.png}{Confirmación de llegada: los integrantes del grupo, con nombre y RUT.}

\begin{nota}
Ya no hace falta escribir a mano quién confirmó la llegada o la
devolución --- el sistema registra automáticamente el nombre de la
persona de personal que tiene la sesión abierta.
\end{nota}

\section{Escanear código de barras (Horizon)}

Las logias también se pueden entregar y devolver escaneando el código de
barras físico de la sala (compatible con el sistema Horizon). El primer
escaneo marca la entrega; el segundo, sobre la misma reserva, marca la
devolución.

\section{Devolución manual de llave}

Desde el detalle de una reserva activa (clic sobre el bloque ocupado) se
puede confirmar la devolución de la llave sin necesidad de escanear un
código de barras.

Si la llave vuelve más de 5 minutos después de la hora de término, al
registrar la devolución (con el botón o escaneando la logia) aparece el
aviso \textbf{Devolución con atraso}: quién está a cargo de la reserva,
cuántos minutos se atrasó y el monto de la multa. \textbf{Sí, generar multa}
la deja pendiente a su nombre (como las de libros: mientras no la pague, no
puede llevarse libros ni recursos de apoyo, ni reservar libros);
\textbf{No, dejarla pasar} cierra el
aviso sin multar. Los montos todavía son provisionales. El escaneo también
sirve pasada la hora: si la llave de un grupo sigue afuera, el escaneo
registra su devolución antes de entregársela al grupo siguiente.

% ======================================================================
\chapter{Reportes}
% ======================================================================

El módulo de Reportes entrega estadísticas del funcionamiento de la
biblioteca para un período concreto: un día, una semana, un mes, un
semestre o un año. Al abrirlo muestra el mes en curso.

\requierepermiso{Reportes}

\captura{07-reportes-prestamos.png}{Reporte de préstamos de un mes, en la vista Gráfico.}

\section{Elegir qué ver}

\begin{itemize}
  \item \textbf{Pestañas:} Préstamos, Ingresos, Uso de salas y Libros más
  prestados.
  \item \textbf{Período:} Día, Semana, Mes, Semestre (enero a junio o julio
  a diciembre) o Año. Las flechas junto al nombre del período pasan al
  anterior o al siguiente; no se puede avanzar más allá del actual.
  \item \textbf{Filtros:} carrera, año de ingreso, sexo y tipo de usuario.
  En Préstamos, además, el tipo de préstamo: bibliográfico (libros) o
  tecnológico (notebooks, audífonos y cargadores).
  \item \textbf{Gráfico o Tabla:} dos formas de ver los mismos datos.
\end{itemize}

\section{Indicadores}

Tres tarjetas resumen el período: el total, comparado con el período
anterior; la carrera con más registros (en Uso de salas, la sala más
usada; en Libros más prestados, el libro), y el horario de mayor uso ---las
dos horas seguidas con más registros---.

\section{Vista Gráfico}

Muestra el total dividido según el período: el día por hora, la semana por
día (los fines de semana, en tono claro, y los días sin atención,
rayados), el mes por semana ---de lunes a domingo, recortadas al mes--- y
el semestre y el año por mes y los
desgloses por carrera, tipo de usuario, sexo, año de ingreso y hora del
día; en este último, las horas de mayor uso van en color fuerte.
Préstamos agrega el desglose \textbf{por tipo de préstamo} ---bibliográfico
o tecnológico, y cuántos de cada recurso---. Uso de
salas agrega las salas más usadas y un mapa de calor con las horas más
pedidas en cada sala, y Libros más prestados, el ranking de títulos.

Las horas siguen el \textbf{horario de atención} de la biblioteca: de
lunes a viernes de 08:45 a 20:30, el sábado de 09:00 a 14:00 y el domingo
cerrado. En el reporte de un día se ven solo las horas de ese día; en un
domingo, el reporte avisa que la biblioteca no atiende.

\captura{08-reportes-logias.png}{Pestaña Uso de salas, con las salas más usadas y el mapa de calor de horas por sala.}

\captura{09-reportes-libros.png}{Pestaña Libros más prestados.}

\section{Vista Tabla}

A la izquierda se elige cómo \textbf{agrupar}: carrera, tipo de usuario,
sexo, año de ingreso, hora del día, día, semana o mes (según el
período) y, en sus pestañas, tipo
de préstamo y recurso, sala o libro. La tabla muestra cada grupo con su cantidad y su porcentaje del
total, y se ordena al hacer clic en los encabezados. \textbf{Exportar
esta tabla} la descarga en Excel, después de confirmar.

\captura{07b-reportes-tabla.png}{Vista Tabla, agrupada por carrera.}

\section{Exportar datos}

\textbf{Exportar} descarga el reporte completo en Excel o CSV, después de
confirmar: el resumen del período, todas las agrupaciones y, en Uso de
salas, las horas por sala. Respeta el período y los filtros elegidos en
pantalla.

% ======================================================================
\chapter{Módulos de Gestión}
\label{cap:gestion}
% ======================================================================

Los módulos secundarios del sistema están agrupados bajo el botón
\textbf{Gestión}, disponible en la barra superior. Al hacer
clic se abre un panel con los módulos ordenados en cuatro secciones:
\textbf{Atención de público}, \textbf{Colección}, \textbf{Consultas e
historiales} y \textbf{Administración del sistema}.

\captura{25-menu-gestiones-admin.png}{Menú de Gestión de una cuenta con rol Administrador/a: todas las secciones y módulos.}

El menú muestra \textbf{solo los módulos que tu rol puede usar}. Por
ejemplo, una cuenta de Atención de público ve un menú mucho más corto:

\captura{25b-menu-asistente.png}{El mismo menú visto por una cuenta con rol Atención de público.}

\section{Usuarios}

\requierepermiso{Consultar usuarios}

Listado y ficha de todos los usuarios institucionales (estudiantes,
docentes y funcionarios), con búsqueda por nombre, apellido, RUT o
carrera, y filtros por tipo de usuario y estado (activo/inactivo).
Como son miles de personas, el listado se muestra de a 25, con los números
de página al pie (por ejemplo, \emph{Mostrando 1 a 25 de 4.650
usuarios}); al buscar o cambiar un filtro vuelve a la primera página.

\captura{10-usuarios.png}{Listado de Usuarios, con filtros de búsqueda.}

Desde cada fila se puede:

\begin{itemize}
  \item \textbf{Habilitar} (solo en quien está inactivo, por ejemplo
  porque su registro venció en el sistema anterior): pide confirmar y desde
  ese momento la persona puede marcar su entrada, pedir préstamos y
  reservas, y entrar al portal. Si tiene multas pendientes, los préstamos y
  las reservas siguen bloqueados hasta que pague. Lo puede hacer todo el
  personal.
  \item \textbf{Ver préstamos:} abre la pantalla de Préstamos con esa
  persona ya buscada.
  \item Generar su \textbf{Constancia de No Multa} en PDF (se bloquea con
  un aviso si el usuario tiene multas pendientes).
  \item \textbf{Desactivar} (solo el Administrador/a): pide confirmar; la
  persona deja de poder marcar su entrada, pedir préstamos y entrar al
  portal, y se cierran las sesiones que tenga abiertas en el portal. Ningún
  otro rol puede recibir este permiso.
\end{itemize}

\section{Multas}
\label{cap:multas}

\requierepermiso{Ver multas}

Reúne las deudas de los usuarios con la biblioteca: pendientes, pagadas y
anuladas. Arriba muestra un resumen (usuarios con deuda, multas
pendientes y monto total pendiente) y dos pestañas:

\begin{itemize}
  \item \textbf{Deuda por usuario:} cada persona con multas pendientes,
  con la cantidad y el monto; \textbf{Ver detalle} muestra sus multas.
  \item \textbf{Todas las multas:} con filtros por estado (pendientes,
  pagadas, anuladas), por motivo (atraso, deterioro, pérdida u otro) y
  búsqueda por RUT o nombre.
\end{itemize}

Las dos pestañas se muestran de a 25 filas, con los números de página al
pie; en \textbf{Todas las multas} las más recientes van primero.

\captura{22-multas.png}{Módulo de Multas, con el resumen, la deuda por usuario y el listado completo.}

Según el rol, cada cuenta puede además:

\begin{itemize}
  \item \textbf{Marcar pagada} (permiso \textit{Registrar pagos}): cuando
  el usuario paga. Queda registrado quién lo hizo y cuándo.
  \item \textbf{Anular} (permiso \textit{Anular multas}): deja sin efecto
  una multa cursada por error. Hay que indicar el motivo; la multa
  \textbf{no se borra}, queda en el historial con el motivo y el nombre de
  quien la anuló.
  \item \textbf{Cursar una multa} (permiso \textit{Cursar multas
  manuales}): para deterioro, pérdida de material u otro motivo, indicando
  el RUT, el motivo, el monto y una descripción. Antes de cursarla, el
  sistema muestra el monto y la persona para confirmar. Las multas por
  atraso no se cursan a mano: se generan solas al devolver.
\end{itemize}

\begin{nota}
Mientras tenga multas pendientes, la persona no puede llevarse libros ni
equipos, ni reservar libros (en mesón ni desde el portal), ni obtener la
Constancia de No Multa. Al pagar o anular sus multas se desbloquea al
instante.
\end{nota}

\section{Listado de Libros}

\requierepermiso{Listados de préstamos y libros}

Catálogo completo de libros (obras). Arriba a la derecha están las
pestañas del módulo ---\textbf{Catalogar}, \textbf{Listado},
\textbf{Historial} y \textbf{Estado de ejemplares}--- para pasar de uno a
otro sin volver al menú; cada una aparece solo si tu rol puede abrirla.

Cada fila muestra el título y el autor, la clasificación Dewey, la
categoría, el tipo de material y una barrita por copia: verde si está
disponible, ámbar si está prestada, violeta si está reservada y gris si
no está en estante (en inventario, en procesos técnicos, dada de baja,
etc.). Debajo se resume cuántas quedan, por ejemplo ``2 de 3
disponibles''.

El buscador encuentra por título, autor, ISBN, clasificación Dewey o
código de barras, sin importar las tildes. Las píldoras de al lado
filtran por \textbf{Estado} de las copias, \textbf{Tipo} de material y
\textbf{Categoría}, y \textbf{Ordenar} cambia el orden entre título y
tipo de recurso.

\captura{12-listado-libros.png}{Listado de Libros del catálogo.}

La flecha a la izquierda de cada fila la despliega: a un lado, la
\textbf{ficha bibliográfica completa} ---con la etiqueta MARC de cada
campo y los de uso interno marcados--- y al otro, cada copia física
(ejemplar) con su código de barras, su ubicación y si está
\emph{Disponible}, \emph{Prestado} o \emph{Reservado} (o en qué estado de
proceso, si no está en estante).

\captura{13-listado-libros-detalle.png}{Detalle desplegado de un título: ficha bibliográfica completa y copias.}

Con el permiso \textit{Catalogar libros}, \textbf{Editar} abre el título
en Catalogación para corregirlo, después de confirmar (los cambios se
guardan recién al pulsar \textbf{Guardar cambios}), y \textbf{+ Agregar
copia} lleva a Catalogación con la obra ya elegida para sumarle un
ejemplar. El botón de tres puntos reúne los atajos a lo que se hace con
ese título en otras pantallas: \textbf{Agregar copia}, \textbf{Historial
de préstamos} y \textbf{Estado de sus ejemplares}, cada uno visible solo
si tu rol tiene el permiso correspondiente.

\section{Catalogación de Libros}

\requierepermiso{Catalogar libros}

La pantalla \textbf{Catalogar material} sirve para ingresar títulos
nuevos al catálogo y para agregar copias a uno ya existente: bajo el
título se elige \textbf{Libro nuevo} o \textbf{Agregar copia a libro
existente}, y las pestañas de arriba a la derecha llevan al
\textbf{Listado}, al \textbf{Historial} y al \textbf{Estado de
ejemplares}.

Lo primero es elegir el \textbf{tipo de material}: define la
\textbf{ficha} del título, es decir, qué etiquetas MARC se catalogan y
cuáles ve el usuario en el portal, según el documento de etiquetas de la
biblioteca. Cada campo indica su etiqueta MARC, y el formulario se divide
en cinco secciones numeradas, cada una de un color:

\begin{itemize}
  \item \textbf{Identificación:} tipo de material, código de barras de la
  primera copia (el botón de al lado sugiere el siguiente disponible), ISBN
  o ISSN, clasificación e idioma.
  \item \textbf{Título y autoría:} título, edición y autores. El primer
  autor que se escribe es el \textbf{autor principal} (100) y el resto
  quedan como \textbf{autores secundarios} (700). El botón
  \textbf{+ Título alternativo, entidad\ldots} muestra los campos
  opcionales (246, 110, 111 y 710).
  \item \textbf{Publicación:} lugar, editorial y año (260; 264 en los
  juegos), descripción física y, según la ficha, serie, volumen y
  frecuencia, o nota de tesis.
  \item \textbf{Materias:} las materias (650), con sugerencias de las que
  ya existen en el catálogo; \textbf{+ Materia de persona, lugar o enlace}
  muestra 600, 651 y 856. En los juegos educativos esta sección se llama
  \textbf{Contenido} (resumen y audiencia).
  \item \textbf{Ejemplar:} ubicación, colección, volumen, carreras
  asociadas y precio de la primera copia, y las notas internas. Nada de
  esto se muestra al usuario.
\end{itemize}

A la izquierda, el índice de \textbf{Secciones} marca en dorado la que
estás viendo y en verde las que ya recorriste (en naranja, si les falta
algo obligatorio); al pulsar una, la página baja hasta ella. A la derecha,
\textbf{Así se verá en el portal} muestra la tarjeta del catálogo mientras
escribes, y \textbf{Obligatorios}, lo mínimo para catalogar: código de
barras, título y autor principal (o la entidad, 110). La creación del
registro (008) y el catalogador (910) se registran solos.

\textbf{Guardar borrador} deja lo escrito guardado en ese computador,
después de confirmar (avisa si reemplaza un borrador anterior): al volver a
la pantalla se ofrece \textbf{Recuperar} o \textbf{Descartar}. Abajo,
\textbf{Títulos catalogados} lista lo ya ingresado, y \textbf{Editar}
carga la ficha arriba para corregirla, también después de confirmar: si
había algo escrito sin guardar, el aviso dice que se perderá. Los datos de
cada copia, como la ubicación o el precio, no se cambian ahí.

\captura{14-catalogacion-libro-nuevo.png}{Formulario de catalogación con la ficha de libro.}

\captura{14b-catalogacion-tesis.png}{Al elegir el tipo Tesis, el formulario cambia a la ficha de trabajo de titulación.}

\begin{center}
\small
\begin{tabularx}{\textwidth}{@{}>{\raggedright\arraybackslash}p{0.2\textwidth}>{\raggedright\arraybackslash}X>{\raggedright\arraybackslash}p{0.27\textwidth}@{}}
\toprule
\textbf{Ficha} & \textbf{Se muestra al usuario} & \textbf{Uso interno} \\
\midrule
Libro & 020, 082, 100, 110, 111, 245, 246, 250, 260, 300, 600, 650, 651, 700, 710, 856 & 008, 041, 490, 500, 504, 910 \\
Trabajo de titulación & 099, 100, 245, 246, 260, 300, 502, 650, 700 (autores secundarios y profesor guía), 710, 856 & 008, 041, 504, 910 \\
Juego educativo & 099, 245, 246, 264, 300, 520, 521 & 008, 041, 910 \\
Publicación seriada (revista) & 022, 099, 245, 246, 260, 300, 362, 650, 856 & 008, 041, 310, 500, 546, 590, 910 \\
\bottomrule
\end{tabularx}
\end{center}

\begin{nota}
Qué ficha usa cada tipo de material se elige en
\textbf{Administración} (sección~\ref{sec:administracion}). La etiqueta
000 (líder) no se registra: se genera al exportar registros en formato
MARC.
\end{nota}

\captura{15-catalogacion-agregar-copia.png}{Agregar una copia: se busca la obra y se completan los datos de la copia nueva.}

Un libro recién catalogado \textbf{no queda disponible para préstamo de
inmediato ni aparece en el portal}: nace en estado ``Inventario'' y debe
pasar por el módulo \textbf{Estado de ejemplares} para marcarse ``En
estante'' antes de poder prestarse.

El botón de generar código de barras sugiere automáticamente el
siguiente código disponible, pero siempre se puede editar antes de
guardar.

\section{Estado de ejemplares}

\requierepermiso{Cambiar estado de un libro}

El estado de proceso es de cada copia física, no de la obra, así que esta
pantalla trabaja copia por copia. A la izquierda, la lista de copias se
busca por título o código de barras ---al escanear el código con el
lector, la copia queda elegida sola--- y se filtra con los botones de
estado, que muestran cuántas copias hay en cada uno. Las copias prestadas
o reservadas llevan además la marca \emph{Prestada} o \emph{Reservada}.

\captura{16-estado-libro.png}{Estado de ejemplares: la lista de copias y, a la derecha, la copia elegida con su flujo de proceso.}

A la derecha se ve la copia elegida con su \textbf{estado actual} y las
opciones para cambiarlo:

\begin{itemize}
  \item \textbf{Flujo de proceso:} el camino normal de una copia nueva,
  \emph{Inventario} $\rightarrow$ \emph{Procesos técnicos} $\rightarrow$
  \emph{Por colocar} $\rightarrow$ \emph{En estante}. Los pasos ya
  recorridos llevan un visto y el actual aparece resaltado.
  \item \textbf{Otras ubicaciones:} \emph{Estantería auxiliar},
  \emph{Colección móvil} o un estado \emph{Personalizado} de los que se
  crean en Administración (al elegirlo se indica cuál).
  \item \textbf{Dar de baja:} retira la copia de la colección. Pide un
  \textbf{motivo obligatorio}, que queda en el historial, y no se puede
  mientras la copia esté prestada o reservada.
\end{itemize}

Elegido el nuevo estado, se puede escribir una nota opcional (por
ejemplo, ``se plastificó la portada'') y se pulsa \textbf{Cambiar
a\ldots}; antes de guardar, el sistema muestra el estado de origen y el
de destino para confirmar. Con el permiso \textit{Historiales}, debajo
aparece el \textbf{Historial de cambios} de esa copia: la fecha, quién
hizo cada cambio y la nota o el motivo.

\begin{importante}
Un ejemplar solo puede prestarse o reservarse si su estado es
\textbf{``En estante''} y además está disponible (no prestado ni
reservado por otra persona). Ambas condiciones son independientes entre
sí.
\end{importante}

\section{Cambio Masivo de Estado}

\requierepermiso{Cambio masivo de estado}

Cambia el estado de muchas copias a la vez, filtrando por ubicación,
estado actual, tipo de material y/o categoría. Requiere elegir al menos
un filtro (no se puede aplicar un cambio masivo sin ningún criterio, para
evitar afectar accidentalmente todo el catálogo). Su encabezado es de
color naranjo, a diferencia del resto, para recordar que es una operación
de riesgo.

\captura{17-cambio-masivo-estado.png}{Vista previa de un cambio masivo de estado, antes de confirmarlo.}

El proceso tiene dos pasos: primero se genera una \textbf{vista previa}
(cuenta cuántos ejemplares coinciden y muestra una muestra de 10), y
luego, para ejecutar el cambio de verdad, el sistema pide
\textbf{escribir la palabra ``CONFIRMAR''} en un cuadro de diálogo --- una
medida de seguridad adicional para una operación que puede afectar
muchas copias a la vez.

Igual que en el cambio de una sola copia, \textbf{dar de baja} en forma
masiva pide un \textbf{motivo obligatorio}, que queda en el historial de
cada copia, y deja fuera las copias que están prestadas o reservadas.

\section{Etiquetas de Código}

\requierepermiso{Catalogar libros}

Genera etiquetas con código de barras real (Code128) a partir de una
lista de códigos --- útil para reimprimir una etiqueta perdida o preparar
una tanda de códigos nuevos antes de catalogar. Sirve para libros,
equipos o cualquier código ya asignado.

\captura{33-etiquetas.png}{Etiquetas de Código, con un rango de códigos listo para imprimir.}

\begin{itemize}
  \item \textbf{Sugerir siguiente (libros):} completa el rango a partir
  del primer código de libro que todavía no se usa.
  \item \textbf{Generar lista con este rango:} arma la lista con el
  prefijo, el número inicial, la cantidad y los dígitos indicados. También
  se pueden pegar códigos a mano, uno por línea.
  \item \textbf{Generar e imprimir:} abre una pestaña nueva con las
  etiquetas y su botón \textbf{Imprimir}.
\end{itemize}

\begin{nota}
Imprimir etiquetas no registra nada en el sistema: los códigos quedan
asignados recién al catalogar cada copia.
\end{nota}

\section{Equipos}

\requierepermiso{Administrar equipos}

Catálogo de audífonos, notebooks y cargadores disponibles para préstamo,
cada uno con su nombre legible, código de barras (lo que se escanea al
prestar) y estado (disponible/activo).

\captura{21-equipos.png}{Catálogo de Equipos.}

Un equipo dado de baja deja de estar disponible para préstamo aunque no
esté prestado en ese momento; y no se puede dar de baja un equipo que
está actualmente prestado. Registrar un equipo y darlo de baja (o
reactivarlo) piden confirmación. Prestar y recibir equipos no requiere
este permiso: se hace desde Préstamos.

\section{Listado de Préstamos}

\requierepermiso{Listados de préstamos y libros}

Listado de consulta de todos los préstamos de libros y equipos, de todos
los usuarios. Las pestañas de arriba muestran cuántos hay en cada grupo:
\textbf{Activos} (aún no devueltos), \textbf{Vencidos} (pasada la fecha de
devolución), \textbf{Devueltos} y \textbf{Todos}.

\captura{11-listado-prestamos.png}{Listado de Préstamos: pestañas por estado, búsqueda y filtro por tipo de ítem.}

El buscador encuentra por RUT (con o sin puntos), nombre, título o código
del ítem, sin importar las tildes, y el selector de la derecha deja ver
solo \textbf{Libros} o solo \textbf{Equipos}. Cada fila indica quién tiene
el ítem, cuándo se prestó (y quién lo prestó) y, en
\textbf{Devolución}, la fecha y cuánto falta (``En 6 días'', ``Vence
hoy'', ``Vencido hace 2 días''), o \emph{Pendiente} en los equipos, que
no tienen fecha fija y se devuelven al terminar la estadía. Si el préstamo
generó una multa, aparece bajo el
estado: pendiente con su monto, o pagada.

\begin{nota}
Este listado es solo de consulta: las devoluciones se registran en
\textbf{Préstamos} (F2) y los pagos de multas, en Préstamos o en el
módulo \textbf{Multas}.
\end{nota}

\section{Historial de Estado}

\requierepermiso{Historiales}

Registro completo de todos los cambios de estado realizados sobre
ejemplares, individuales o masivos, con filtros por título/código,
identificador de lote (para ver qué se cambió junto en una operación
masiva), funcionario que hizo el cambio, y rango de fechas.

\captura{18-historial-estado.png}{Historial de Estado de ejemplares.}

\section{Historial de Libros}

\requierepermiso{Historiales}

A diferencia del anterior (que es sobre \emph{estado}), este módulo es
sobre \emph{préstamos}. A la izquierda están los libros de la \textbf{A a
la Z}, cada uno con su cantidad de préstamos, y se buscan por título, ISBN
o código de barras.

\captura{19-historial-prestamos-libro.png}{Historial de libros: los libros con su cantidad de préstamos y, a la derecha, el libro elegido con sus indicadores y cada préstamo.}

Al elegir un libro se ven cuatro indicadores ---\textbf{Préstamos},
\textbf{Duración promedio} de los ya devueltos, \textbf{Devueltos con
atraso} y la fecha del \textbf{Último préstamo}--- y debajo cada préstamo
con la copia, el usuario, las fechas de préstamo y devolución, los días y
su resultado: \emph{A tiempo}, \emph{Atraso} (con los días), \emph{En
curso} o \emph{Vencido}. \textbf{Ver ficha} abre la ficha bibliográfica
completa del título. Desde el Listado de Libros, el atajo
\textbf{Historial de préstamos} abre esta pantalla con ese libro ya
elegido.

\section{Registro Absoluto}
\label{cap:registro-absoluto}

\requierepermiso{Registro absoluto}

Un registro unificado de \textbf{todas las operaciones que se solicitan
en el sistema}: préstamos, reservas de sala, reservas de libro y
entradas, todas juntas en una sola tabla, ordenadas por fecha.

\captura{24-registro-absoluto.png}{Registro Absoluto, con filtros por fecha, tipo de operación y búsqueda.}

Es la herramienta pensada para auditorías y consultas históricas que
ningún otro módulo cubre por sí solo --- por ejemplo, ver todo lo que
hizo una persona puntual (por RUT o nombre) a través de varios meses, sin
importar si fue un préstamo, una reserva de sala o una simple entrada.

\begin{itemize}
  \item \textbf{Rango de fechas:} por defecto muestra los últimos 30
  días; se puede ajustar libremente.
  \item \textbf{Tipo de operación:} casillas para incluir o excluir
  préstamos, reservas de sala, reservas de libro y/o entradas.
  \item \textbf{Búsqueda:} por nombre o RUT del usuario involucrado.
  \item \textbf{Atendido por:} muestra el nombre del funcionario que
  procesó la operación, cuando corresponde. Queda vacío cuando la
  operación la realizó el propio usuario desde el portal de autoservicio
  --- eso no es un dato faltante, es información real (nadie de personal
  intervino en esa operación puntual).
  \item \textbf{Exportar a Excel:} descarga exactamente lo que está
  filtrado en pantalla en ese momento, después de confirmar el rango de
  fechas y la cantidad de operaciones. El archivo incluye nombres y RUT:
  guárdalo solo donde corresponda.
\end{itemize}

\section{Código QR de Acceso}

\requierepermiso{Mostrar el código QR de acceso}

Genera y muestra el código QR único que los usuarios pueden escanear
desde el portal para marcar su entrada sin escribir su RUT.

\captura{20-codigo-qr.png}{Código QR de acceso, con opción de descarga e impresión.}

El texto de 16 letras y números bajo el QR es el mismo código escrito: el
QR no contiene nada más. Es igual para todos y no identifica a nadie;
sirve para demostrar que la persona está en la biblioteca. La identidad
la aporta la sesión del portal: quien escanea ya entró con su RUT y
contraseña, así que la entrada queda a su nombre, con vía \textbf{QR
Móvil}. El código no vence solo; si se sospecha que circula una foto de
él fuera de la biblioteca, conviene regenerarlo.

\begin{importante}
Al \textbf{regenerar} el código (permiso \textit{Regenerar el código
QR}), el anterior deja de funcionar de inmediato --- cualquier copia
impresa o pantalla con el QR viejo debe reemplazarse. Sin ese permiso, la
pantalla solo muestra, descarga e imprime el código vigente.
\end{importante}

\section{Administración}
\label{sec:administracion}

Configuración general del sistema, catálogos administrables y estudiantes,
en tres pestañas. Cada una aparece solo si el rol tiene su permiso.

\captura{23-administracion.png}{Página de Administración, pestaña de configuración institucional.}

\begin{itemize}
  \item \textbf{Configuración institucional} (permiso
  \textit{Configuración institucional}): nombre y cargo de quien firma la
  Constancia de No Multa.
  \item \textbf{Catálogos} (permiso \textit{Catálogos del sistema}):
  \begin{itemize}
    \item \textbf{Estados personalizados de libro:} estados adicionales,
    más allá de los siete fijos del sistema, para necesidades específicas
    de la biblioteca (por ejemplo, ``Restauración''). Se pueden desactivar
    sin perder el historial de los ejemplares que ya los tuvieron
    asignados.
    \item \textbf{Ubicaciones:} sedes o ubicaciones físicas donde puede
    estar un ejemplar.
    \item \textbf{Tipos de material:} por ejemplo Libro, Revista, Tesis o DVD.
    Cada tipo elige su \textbf{ficha} (Libro, Trabajo de titulación,
    Juego educativo o Publicación seriada), que define qué etiquetas se
    catalogan y cuáles ve el usuario (ver Catalogación de Libros).
    Cambiar la ficha de un tipo existente pide confirmación; lo ya
    catalogado conserva sus datos.
  \end{itemize}
  \item \textbf{Estudiantes} (permiso \textit{Crear y editar estudiantes}):
  registrar estudiantes y otros usuarios, corregir sus datos con
  \textbf{Editar} y cargarlos en forma masiva desde una planilla Excel o
  CSV, también el reporte de matrícula de la universidad tal como viene. La
  \textbf{carga masiva} primero muestra qué pasará con cada fila, sin
  guardar nada:
  \begin{itemize}
    \item \textbf{Se cargará:} la persona no está en el sistema y se crea,
    activa.
    \item \textbf{Se actualizará:} ya está (mismo RUT o documento). Sus
    datos se completan o corrigen con los de la planilla (nombre,
    apellidos, carrera, año de ingreso, sexo, correo) y, si estaba
    inactiva, queda activa. Se ve cada cambio como ``antes $\rightarrow$
    después''. Lo que en la planilla viene vacío no borra nada, y el tipo
    (estudiante, docente o funcionario) solo cambia si la planilla trae esa
    columna.
    \item \textbf{Sin cambios:} ya está y sus datos coinciden con la
    planilla.
    \item \textbf{Repetido en el archivo} o \textbf{Con errores:} no se
    carga; el detalle dice por qué y se puede descargar para corregir la
    planilla.
  \end{itemize}
  Nada se guarda hasta marcar la casilla de confirmación y pulsar el botón.
  Cada persona actualizada queda en el Registro Absoluto con lo que cambió.
\end{itemize}

\begin{nota}
Si el correo de la planilla ya lo usa otra persona, ese correo no se
cambia y el resto de los datos sí. Y si el número de un RUT coincide con
el de un DNI o pasaporte extranjero de la planilla, la fila no se carga
para no mezclar a dos personas: hay que revisarla a mano.
\end{nota}

\captura{23b-administracion-tipos.png}{Tipos de material, cada uno con la ficha que usa.}

\begin{nota}
Los códigos de barra para imprimir, que antes estaban en esta pantalla,
ahora son el módulo \textbf{Etiquetas de Código}.
\end{nota}

\section{Personal y Roles}
\label{cap:personal-roles}

\requierepermiso{Cuentas y roles del personal}

Administra las cuentas del personal de la biblioteca y qué puede hacer
cada rol. Tiene dos pestañas.

\subsection{Cuentas del personal}

Lista a cada persona con su correo, rol, estado y último acceso.

\captura{34-personal-cuentas.png}{Personal y Roles, pestaña de cuentas del personal.}

\begin{itemize}
  \item \textbf{Nueva cuenta:} nombre completo, correo, rol y una
  contraseña inicial (el formulario propone una generada al azar, con un botón para copiarla). La persona
  entra con su correo y esa contraseña, y después puede cambiarla desde
  el menú de su cuenta (su nombre, arriba a la derecha).
  \item \textbf{Editar:} cambia nombre, correo o rol.
  \item \textbf{Restablecer contraseña:} asigna una nueva a quien la
  olvidó: se escribe o se genera al azar con \textbf{Generar}. La actual no
  se puede ver, porque el sistema la guarda cifrada. Sus sesiones abiertas
  se cierran.
  \item \textbf{Desactivar:} para quien deja de trabajar en la
  biblioteca. Su sesión se cierra al instante y su historial (préstamos,
  entradas, cambios) se conserva. Se puede reactivar después.
\end{itemize}

Todas estas acciones piden confirmación. Al editar se muestra lo que
cambia (por ejemplo, el rol anterior tachado y el nuevo, con un aviso si
el nuevo tiene acceso total), y al restablecer, la contraseña nueva para
anotarla antes de confirmar.

\subsection{Roles y permisos}

Muestra cada rol con su descripción y un cuadro \textbf{¿Qué puede hacer
cada rol?} que compara los permisos de todos los roles.

\captura{35-personal-roles.png}{Pestaña de roles, con la comparación de permisos entre roles.}

\begin{itemize}
  \item \textbf{Editar permisos:} marca o desmarca lo que incluye un rol.
  Antes de guardar, el sistema muestra qué permisos se agregan y cuáles
  se quitan, y a cuántas cuentas afecta, y pide confirmar. El cambio
  aplica de inmediato a todas las cuentas con ese rol.
  \item \textbf{Nuevo rol:} por ejemplo ``Ayudante de sala'', que solo
  registra entradas. Se puede partir desde los permisos de un rol
  existente. También pide confirmación, con un aviso si no tiene ningún
  permiso marcado.
  \item \textbf{Eliminar un rol:} solo si ninguna cuenta lo tiene
  asignado.
\end{itemize}

\begin{importante}
Solo el rol con acceso total (Administrador/a) puede gestionar cuentas y
roles, para que nadie pueda darse más permisos a sí mismo. Por la misma
razón, el sistema no permite desactivar la propia cuenta, ni dejar la
biblioteca sin ninguna cuenta activa con acceso total.
\end{importante}

% ======================================================================
\chapter{Portal del Usuario}
% ======================================================================

El Portal del Usuario es un sitio de autoservicio independiente, pensado
para que estudiantes, docentes y funcionarios resuelvan trámites simples
por su cuenta, sin pasar por mesón.

El portal se adapta al equipo desde el que se abre. En un computador
(ventanas de 1024 píxeles de ancho o más) muestra un panel de inicio con
la actividad de la persona, las dos formas de marcar entrada lado a lado
y el catálogo junto a \textbf{Mis reservas}. En teléfonos y tablets usa
una sola columna, con botones grandes pensados para el tacto.

En la barra superior, las pestañas \textbf{Inicio}, \textbf{Entrada},
\textbf{Catálogo} y \textbf{Salas} llevan a cada sección. A la derecha
aparecen el nombre y el tipo de usuario; al pulsarlos se abre un menú con
el RUT de la cuenta, \textbf{Mi perfil}, \textbf{Mis ingresos} y
\textbf{Cerrar sesión}. En el teléfono, una barra fija abajo lleva a
Inicio, Catálogo, Entrada (el botón central), Salas y Perfil.

Las capturas de este capítulo muestran el formato de computador; el de
teléfono se ve así:

\captura{27b-portal-telefono.png}{El portal en un teléfono: inicio, marcar entrada, catálogo y salas (la línea de tiempo se desliza de lado).}

\section{Inicio de sesión}

Se accede con el RUT y la contraseña personal del usuario (distinta de
las credenciales de personal). Los estudiantes extranjeros marcan
\textbf{Pasaporte / DNI} en el selector sobre la casilla y escriben su
número; el navegador recuerda esa elección para la próxima vez.

\captura{26-portal-login.png}{Inicio de sesión del Portal del Usuario.}

El campo de RUT se autocompleta con los puntos y el guión a medida que
se escribe --- no hace falta tipearlos a mano.

\captura{26-portal-login-completo.png}{RUT y contraseña completados, listos para iniciar sesión.}

\begin{nota}
Si un usuario olvida su contraseña, el sistema no ofrece recuperación en
línea --- debe acercarse al mesón de la biblioteca con su credencial para
que el personal se la restablezca.
\end{nota}

\section{Página de Inicio}

En el computador, el inicio es un panel con lo más importante del día:

\captura{27-portal-home.png}{Inicio del portal en el computador: ingreso de hoy, actividad y accesos rápidos.}

\begin{itemize}
  \item \textbf{¿Llegaste a la biblioteca?:} indica si la persona ya marcó
  su ingreso hoy (y a qué hora) y ofrece \textbf{Escanear QR} o
  \textbf{Usar mi RUT}, que llevan directo a Marcar Entrada con la cámara
  encendida o con el cursor en el campo del RUT.
  \item \textbf{Ingresos de hoy:} cuántas personas marcaron su entrada
  hoy, frente al aforo de la biblioteca. Como el sistema no registra
  salidas, no es la cantidad de personas que están adentro en este
  momento.
  \item \textbf{Buscador:} busca en el catálogo por título, autor o ISBN.
  \item \textbf{Tu actividad:} los préstamos vigentes con su fecha de
  devolución (con un aviso si vencen en los próximos días o están
  atrasados), las reservas de libros, las reservas de sala de hoy y el
  estado de multas. Sin multas pendientes, desde ahí se descarga la
  \textbf{Constancia de No Multa} en PDF.
  \item \textbf{Accesos rápidos:} reservar sala, explorar el catálogo y
  \textbf{Recursos digitales} (\url{https://umag.elogim.com/}), la base de
  datos digital de la universidad, que es un sitio externo al sistema de
  biblioteca.
\end{itemize}

Arriba a la derecha, una etiqueta indica si la biblioteca está abierta
(«Abierta · cierra a las 20:30») o, si está cerrada, cuándo vuelve a
abrir (por ejemplo, «Cerrada · abre el lunes a las 08:45»), según el
horario de atención: lunes a viernes de 08:45 a 20:30 y sábado de 09:00 a
14:00.

En el teléfono, el inicio saluda con la fecha, el horario de la
biblioteca y los ingresos del día. Debajo están el aviso
\textbf{¿Llegaste a la biblioteca?} con \textbf{Escanear QR} y
\textbf{Usar mi RUT}, el buscador del catálogo y \textbf{Explorar}:
Catálogo, Reservar sala, Recursos digitales y Mis ingresos. La Constancia
de No Multa, en el teléfono, está en \textbf{Perfil}.

\section{Marcar Entrada}

Permite registrar el propio ingreso de dos formas, cada una en su
tarjeta. Arriba a la derecha se indica si la biblioteca está abierta y
hasta qué hora.

\captura{28-portal-entrada.png}{Marcar Entrada: escanear el código QR o ingresar con el RUT.}

\begin{itemize}
  \item \textbf{Escanear código QR} (recomendado en el celular):
  \textbf{Abrir cámara} enciende la cámara dentro de la misma tarjeta; al
  leer el código QR de la biblioteca, la entrada queda registrada. La
  primera vez, el navegador pide permiso para usar la cámara.
  \item \textbf{Ingresar con RUT} (ideal en el computador): se escribe el
  RUT sin puntos ni guion ---el sistema lo formatea--- y se pulsa
  \textbf{Registrar entrada} o \textbf{Enter}.
\end{itemize}

Si algo falla ---por ejemplo, un RUT que no es el de la cuenta--- el aviso
aparece en la misma tarjeta, sin salir de la pantalla.

\captura{28b-portal-entrada-rut.png}{El RUT escrito no es el de la cuenta: el aviso aparece en la tarjeta.}

\begin{nota}
El sistema exige que el RUT ingresado coincida exactamente con el de la
cuenta que inició sesión --- un usuario no puede marcar la entrada de
otra persona desde su propia sesión del portal.
\end{nota}

Con el RUT correcto, la entrada queda registrada de inmediato y el
sistema muestra una bienvenida antes de volver solo al inicio del
portal.

\captura{28c-portal-entrada-exito.png}{Confirmación de entrada registrada.}

\section{Catálogo}

Consulta de disponibilidad de libros y revistas. El buscador acepta
título, autor, área o ISBN (con o sin guiones) e indica cuántos
resultados hay; los botones \textbf{Todos} y \textbf{Disponibles ahora} y
el selector de \textbf{Área} filtran la lista. Al pulsar \textbf{Área}
aparecen las áreas con más títulos; para encontrar otra, se escriben unas
letras de su nombre y se elige de la lista (la cruz $\times$ quita el filtro).

\captura{29-portal-catalogo.png}{Catálogo del portal, con filtros, disponibilidad y reserva directa.}

Cada título muestra una portada con sus iniciales (del color de su área),
su estado y la acción que corresponde:

\begin{itemize}
  \item \textbf{Disponible}, con cuántos ejemplares hay en estantería:
  botón \textbf{Reservar}.
  \item \textbf{Prestado}, con la fecha aproximada en que vuelve una copia:
  botón \textbf{Avisarme al volver}, que deja a la persona en la lista de
  espera. \textbf{No disponible} indica que no hay copias libres aunque
  ninguna esté prestada (por ejemplo, porque está apartada para otra
  persona).
  \item \textbf{Lista para retirar} o \textbf{En lista de espera}, si la
  persona ya lo reservó: en lugar del botón aparece \textbf{En tus
  reservas}.
\end{itemize}

\textbf{Ver ficha}, junto a cada título, muestra sus datos
bibliográficos públicos según el tipo de material (ISBN, clasificación,
pie de imprenta, materias, acceso electrónico\ldots). En el computador
se despliega dentro de la lista; en el teléfono se abre una ventana con
el resumen del título ---portada, estado, ficha y el botón para
reservar--- sin alargar la página. Los datos de uso interno de la
biblioteca no se muestran.

\captura{29b-portal-ficha.png}{Ficha de un título en el portal: solo los datos públicos.}

Reservar siempre pide confirmación. \textbf{Reservar} y
\textbf{Avisarme al volver} abren una ventana con el título y lo que va a
pasar: en el primer caso se aparta una copia a nombre de la persona para
retirarla en el mesón antes de la fecha límite; en el segundo, queda en la
lista de espera y, cuando llega su turno, se le aparta una copia y se le
avisa por correo si tiene uno registrado. Nada se registra hasta pulsar
\textbf{Confirmar}; \textbf{Cancelar} cierra la ventana sin reservar.

\captura{29c-portal-confirmar-reserva.png}{Confirmación antes de reservar un libro.}

\textbf{Mis reservas} (a la derecha en el computador; sobre la lista en el
teléfono) muestra cada reserva con su estado, el plazo para retirarla o
el lugar en la fila, y permite cancelarla o salir de la fila (también
con confirmación). Después de cancelar o salir de la fila hay que esperar
cinco minutos para pedir otro libro, para que nadie pida y cancele en
seguidilla; mientras tanto, el portal avisa a qué hora se puede volver a
pedir. Debajo, \textbf{¿Cómo retiro
un libro?} recuerda que, cuando la reserva diga «Lista para retirar», hay
que pasar al mesón de préstamo con el RUT dentro del plazo indicado. Si no
se retira a tiempo, la reserva se cancela sola al día siguiente y la copia
pasa a la siguiente persona de la fila.

\begin{nota}
Con multas pendientes, el portal no permite reservar ni unirse a la lista
de espera: muestra el monto adeudado y pide pagarlo en la biblioteca.
\end{nota}

\section{Mi perfil y Mis ingresos}

\textbf{Mi perfil} (en el menú de la cuenta, o \textbf{Perfil} en la
barra inferior del teléfono) muestra el nombre, el tipo de usuario, el
RUT, la carrera y el correo, y reúne la \textbf{Constancia de No Multa}
(se bloquea si hay multas pendientes), el acceso a \textbf{Mis ingresos}
y \textbf{Cerrar sesión}.

\textbf{Mis ingresos} muestra cuántas veces marcó la persona su entrada
este mes y su historial reciente, agrupado por día, con la hora y cómo se
registró: con el código QR, con su RUT en el portal o en el mesón.

\captura{27c-portal-perfil-ingresos.png}{Mi perfil y Mis ingresos, en el teléfono.}

\section{Salas de Estudio}

Reserva de logias desde el portal, con la misma línea de tiempo por sala
que usa el personal: lo que ya pasó aparece rayado y la línea azul, con
la etiqueta \textbf{Ahora}, marca la hora actual. Las reservas de 30
minutos, por ser cortas, se dibujan como un punto; al pasar el cursor
sobre cualquier bloque se ve su horario. El botón \textbf{¿Cómo reservo
una sala?} despliega los pasos, y la página se actualiza sola cada 30
segundos. Lo que queda fuera del horario de la biblioteca también aparece
rayado (el sábado, desde las 14:00, con un aviso), y los domingos, como la
biblioteca no atiende, la pantalla lo avisa y no permite reservar. En el
teléfono, la línea de tiempo se desliza hacia el lado, con
el nombre de cada sala fijo a la izquierda, y se abre mostrando la hora
actual.

\captura{30-portal-salas.png}{Reserva de salas de estudio desde el portal.}

Al hacer clic en un espacio disponible se abre el mismo formulario de
reserva que usa el personal, con una diferencia: el RUT de la propia
persona ya viene completado automáticamente en el primer campo. En el
portal las reservas empiezan a la hora o a la media hora; la opción
\textbf{Reservar ahora}, que empieza en el minuto exacto, es solo del
personal.

\captura{30b-portal-salas-modal-reserva.png}{Formulario de reserva desde el portal --- el primer RUT ya viene precargado con la cuenta propia.}

Las reservas propias aparecen en otro color. Al pulsar una, el portal
pregunta \textbf{¿Cancelar tu reserva?} antes de liberarla. Solo se
cancela una reserva que todavía no se usa: si el grupo ya está en la sala,
se cierra en el mesón al devolver la llave, y si pasaron los 15 minutos
sin que llegaran, queda registrada como inasistencia. Las reservas
terminadas o con inasistencia quedan en el historial y no se pueden
cancelar.

Quien hizo la reserva puede cancelar una al día sin consecuencias; desde
la segunda, tiene que esperar cinco minutos para reservar otra sala, y el
portal se lo advierte antes de confirmar la cancelación. Un integrante
(por ejemplo, alguien a quien incluyeron sin preguntarle) no cancela la
reserva de todo el grupo: al pulsarla, el portal pregunta \textbf{¿Salir
de esta reserva?}, y si confirma, sale de la lista y el resto del grupo la
mantiene. Como las salas se reservan para dos o más personas, si en el
grupo quedaría una sola, la reserva se cancela. Salirse no le pesa a
nadie, y quien está esperando igual puede ir como integrante en las
reservas de otros.

\begin{importante}
Desde el portal, los usuarios \textbf{solo pueden reservar salas para el
día de hoy} --- no existe la opción de reservar con anticipación para
otro día. El personal, en cambio, sí puede reservar para cualquier fecha
desde el panel interno.
\end{importante}

% ======================================================================
\chapter{Preguntas Frecuentes}
% ======================================================================

\begin{description}[leftmargin=0pt,style=nextline]
  \item[¿Por qué no puedo prestar un libro que aparece en el catálogo?]
  Verifica su estado en el módulo \textbf{Estado de ejemplares}: solo los
  ejemplares ``En estante'' y sin prestar/reservar por otra persona se
  pueden prestar. Un libro recién catalogado nace en ``Inventario'', no
  en estante.

  \item[¿Un usuario con multa pendiente puede seguir pidiendo libros?]
  No. Mientras tenga multas pendientes no puede llevarse libros ni
  equipos, ni reservar libros, ni obtener la Constancia de No Multa.
  Apenas se registra el pago (o se anula la multa), puede volver a pedir.

  \item[¿Por qué no encuentro el módulo X en mi menú?]
  Porque tu rol no incluye ese permiso. Los recuadros ``Requiere
  permiso'' de este manual indican cuál hace falta; si lo necesitas para
  tu trabajo, pídele a un administrador que lo agregue a tu rol o que te
  asigne otro (Personal y Roles).

  \item[Olvidé mi contraseña de personal. ¿Qué hago?]
  Un administrador puede restablecerla desde Personal y Roles
  (botón \textbf{Contraseña} en tu fila). Después, cámbiala por una
  propia desde el menú de tu cuenta (tu nombre, arriba a la derecha).

  \item[Un compañero dejó la biblioteca. ¿Borro su cuenta?]
  No: se \textbf{desactiva} desde Personal y Roles. Así no puede volver a
  entrar, y su nombre sigue apareciendo en los préstamos, entradas y
  cambios que atendió.

  \item[Una multa se cursó por error. ¿Se puede borrar?]
  No se borra: se \textbf{anula} desde el módulo Multas, indicando el
  motivo (requiere el permiso \textit{Anular multas}). Queda en el
  historial con el motivo y el nombre de quien la anuló.

  \item[¿Puedo reservar una sala para la próxima semana desde el portal?]
  No. El portal solo permite reservar para el día de hoy (y nunca un
  domingo, que la biblioteca no atiende). El personal sí puede reservar
  para otras fechas desde el panel interno.

  \item[¿Qué pasa si un grupo reserva una sala y no llega?]
  Tiene 15 minutos desde la hora de inicio de su reserva para confirmar
  su llegada. Si no lo hace, la sala se libera automáticamente y queda
  disponible para otra persona.

  \item[¿Cómo sé quién atendió una operación puntual?]
  El módulo \textbf{Registro Absoluto} (permiso \textit{Registro
  absoluto}) muestra el nombre del funcionario que atendió cada préstamo,
  reserva de sala o reserva de libro --- salvo que la haya hecho el propio
  usuario desde el portal, en cuyo caso no hay funcionario involucrado.

  \item[¿El sistema funciona si no hay conexión con el servidor?]
  No: ninguna pantalla muestra datos de ejemplo o ficticios si falla la
  conexión. Si aparece el aviso ``No se pudo conectar con el servidor'',
  ningún dato mostrado en pantalla es real hasta que la conexión se
  restablezca.
\end{description}

% ======================================================================
\chapter*{Glosario}
\addcontentsline{toc}{chapter}{Glosario}
% ======================================================================

\begin{description}[leftmargin=0pt,style=nextline]
  \item[Ejemplar] Una copia física concreta de un libro (con su propio
  código de barras). Un mismo libro (obra) puede tener varios ejemplares.

  \item[Estado de proceso] Situación física/administrativa de un
  ejemplar: en inventario, en procesos técnicos, por colocar, en
  estante, en estantería auxiliar, en colección móvil, dado de baja, o
  un estado personalizado.

  \item[Ficha bibliográfica] Los datos de un título organizados por
  etiquetas MARC. Qué etiquetas lleva, y cuáles ve el usuario, depende del
  tipo de material.

  \item[Disponible] Indica si un ejemplar está prestado/reservado en
  este momento o no. Es independiente del estado de proceso --- un
  ejemplar puede estar ``en estante'' pero no disponible porque otra
  persona lo tiene prestado.

  \item[Logia] Nombre que usa la biblioteca para sus salas de estudio
  individuales (15 en total), reservables en horario continuo por hasta
  dos horas cada vez.

  \item[Multa] Deuda de un usuario con la biblioteca. Las de atraso se
  generan solas al devolver un libro fuera de plazo; las de deterioro,
  pérdida u otro motivo se cursan a mano. Se pagan o se anulan, pero
  nunca se borran.

  \item[Permiso] Una acción concreta del panel de personal (por ejemplo
  ``Catalogar libros'' o ``Anular multas''). Los permisos se agrupan en
  roles.

  \item[Rol] Conjunto de permisos que se asigna a una cuenta del
  personal: Atención de público, Bibliotecario/a, Administrador/a u otro
  creado por la biblioteca.

  \item[Cola de espera] Lista de personas esperando que se libere una
  copia de un libro agotado, en orden de llegada (primero en pedirlo,
  primero en recibirlo cuando se libere una copia).

  \item[Registro Absoluto] Vista unificada de todas las operaciones
  (préstamos, reservas de sala, reservas de libro y entradas) del
  sistema, con filtros de fecha, tipo y búsqueda por persona.

  \item[Portal del Usuario] Sitio de autoservicio separado del panel de
  personal, para que estudiantes, docentes y funcionarios resuelvan
  trámites simples por su cuenta.
\end{description}

\end{document}
